% Options for packages loaded elsewhere
\PassOptionsToPackage{unicode}{hyperref}
\PassOptionsToPackage{hyphens}{url}
\documentclass[
]{article}
\usepackage{xcolor}
\usepackage{amsmath,amssymb}
\setcounter{secnumdepth}{-\maxdimen} % remove section numbering
\usepackage{iftex}
\ifPDFTeX
  \usepackage[T1]{fontenc}
  \usepackage[utf8]{inputenc}
  \usepackage{textcomp} % provide euro and other symbols
\else % if luatex or xetex
  \usepackage{unicode-math} % this also loads fontspec
  \defaultfontfeatures{Scale=MatchLowercase}
  \defaultfontfeatures[\rmfamily]{Ligatures=TeX,Scale=1}
\fi
\usepackage{lmodern}
\ifPDFTeX\else
  % xetex/luatex font selection
\fi
% Use upquote if available, for straight quotes in verbatim environments
\IfFileExists{upquote.sty}{\usepackage{upquote}}{}
\IfFileExists{microtype.sty}{% use microtype if available
  \usepackage[]{microtype}
  \UseMicrotypeSet[protrusion]{basicmath} % disable protrusion for tt fonts
}{}
\makeatletter
\@ifundefined{KOMAClassName}{% if non-KOMA class
  \IfFileExists{parskip.sty}{%
    \usepackage{parskip}
  }{% else
    \setlength{\parindent}{0pt}
    \setlength{\parskip}{6pt plus 2pt minus 1pt}}
}{% if KOMA class
  \KOMAoptions{parskip=half}}
\makeatother
\usepackage{color}
\usepackage{fancyvrb}
\newcommand{\VerbBar}{|}
\newcommand{\VERB}{\Verb[commandchars=\\\{\}]}
\DefineVerbatimEnvironment{Highlighting}{Verbatim}{commandchars=\\\{\}}
% Add ',fontsize=\small' for more characters per line
\newenvironment{Shaded}{}{}
\newcommand{\AlertTok}[1]{\textcolor[rgb]{1.00,0.00,0.00}{\textbf{#1}}}
\newcommand{\AnnotationTok}[1]{\textcolor[rgb]{0.38,0.63,0.69}{\textbf{\textit{#1}}}}
\newcommand{\AttributeTok}[1]{\textcolor[rgb]{0.49,0.56,0.16}{#1}}
\newcommand{\BaseNTok}[1]{\textcolor[rgb]{0.25,0.63,0.44}{#1}}
\newcommand{\BuiltInTok}[1]{\textcolor[rgb]{0.00,0.50,0.00}{#1}}
\newcommand{\CharTok}[1]{\textcolor[rgb]{0.25,0.44,0.63}{#1}}
\newcommand{\CommentTok}[1]{\textcolor[rgb]{0.38,0.63,0.69}{\textit{#1}}}
\newcommand{\CommentVarTok}[1]{\textcolor[rgb]{0.38,0.63,0.69}{\textbf{\textit{#1}}}}
\newcommand{\ConstantTok}[1]{\textcolor[rgb]{0.53,0.00,0.00}{#1}}
\newcommand{\ControlFlowTok}[1]{\textcolor[rgb]{0.00,0.44,0.13}{\textbf{#1}}}
\newcommand{\DataTypeTok}[1]{\textcolor[rgb]{0.56,0.13,0.00}{#1}}
\newcommand{\DecValTok}[1]{\textcolor[rgb]{0.25,0.63,0.44}{#1}}
\newcommand{\DocumentationTok}[1]{\textcolor[rgb]{0.73,0.13,0.13}{\textit{#1}}}
\newcommand{\ErrorTok}[1]{\textcolor[rgb]{1.00,0.00,0.00}{\textbf{#1}}}
\newcommand{\ExtensionTok}[1]{#1}
\newcommand{\FloatTok}[1]{\textcolor[rgb]{0.25,0.63,0.44}{#1}}
\newcommand{\FunctionTok}[1]{\textcolor[rgb]{0.02,0.16,0.49}{#1}}
\newcommand{\ImportTok}[1]{\textcolor[rgb]{0.00,0.50,0.00}{\textbf{#1}}}
\newcommand{\InformationTok}[1]{\textcolor[rgb]{0.38,0.63,0.69}{\textbf{\textit{#1}}}}
\newcommand{\KeywordTok}[1]{\textcolor[rgb]{0.00,0.44,0.13}{\textbf{#1}}}
\newcommand{\NormalTok}[1]{#1}
\newcommand{\OperatorTok}[1]{\textcolor[rgb]{0.40,0.40,0.40}{#1}}
\newcommand{\OtherTok}[1]{\textcolor[rgb]{0.00,0.44,0.13}{#1}}
\newcommand{\PreprocessorTok}[1]{\textcolor[rgb]{0.74,0.48,0.00}{#1}}
\newcommand{\RegionMarkerTok}[1]{#1}
\newcommand{\SpecialCharTok}[1]{\textcolor[rgb]{0.25,0.44,0.63}{#1}}
\newcommand{\SpecialStringTok}[1]{\textcolor[rgb]{0.73,0.40,0.53}{#1}}
\newcommand{\StringTok}[1]{\textcolor[rgb]{0.25,0.44,0.63}{#1}}
\newcommand{\VariableTok}[1]{\textcolor[rgb]{0.10,0.09,0.49}{#1}}
\newcommand{\VerbatimStringTok}[1]{\textcolor[rgb]{0.25,0.44,0.63}{#1}}
\newcommand{\WarningTok}[1]{\textcolor[rgb]{0.38,0.63,0.69}{\textbf{\textit{#1}}}}
\usepackage{graphicx}
\makeatletter
\newsavebox\pandoc@box
\newcommand*\pandocbounded[1]{% scales image to fit in text height/width
  \sbox\pandoc@box{#1}%
  \Gscale@div\@tempa{\textheight}{\dimexpr\ht\pandoc@box+\dp\pandoc@box\relax}%
  \Gscale@div\@tempb{\linewidth}{\wd\pandoc@box}%
  \ifdim\@tempb\p@<\@tempa\p@\let\@tempa\@tempb\fi% select the smaller of both
  \ifdim\@tempa\p@<\p@\scalebox{\@tempa}{\usebox\pandoc@box}%
  \else\usebox{\pandoc@box}%
  \fi%
}
% Set default figure placement to htbp
\def\fps@figure{htbp}
\makeatother
% definitions for citeproc citations
\NewDocumentCommand\citeproctext{}{}
\NewDocumentCommand\citeproc{mm}{%
  \begingroup\def\citeproctext{#2}\cite{#1}\endgroup}
\makeatletter
 % allow citations to break across lines
 \let\@cite@ofmt\@firstofone
 % avoid brackets around text for \cite:
 \def\@biblabel#1{}
 \def\@cite#1#2{{#1\if@tempswa , #2\fi}}
\makeatother
\newlength{\cslhangindent}
\setlength{\cslhangindent}{1.5em}
\newlength{\csllabelwidth}
\setlength{\csllabelwidth}{3em}
\newenvironment{CSLReferences}[2] % #1 hanging-indent, #2 entry-spacing
 {\begin{list}{}{%
  \setlength{\itemindent}{0pt}
  \setlength{\leftmargin}{0pt}
  \setlength{\parsep}{0pt}
  % turn on hanging indent if param 1 is 1
  \ifodd #1
   \setlength{\leftmargin}{\cslhangindent}
   \setlength{\itemindent}{-1\cslhangindent}
  \fi
  % set entry spacing
  \setlength{\itemsep}{#2\baselineskip}}}
 {\end{list}}
\usepackage{calc}
\newcommand{\CSLBlock}[1]{\hfill\break\parbox[t]{\linewidth}{\strut\ignorespaces#1\strut}}
\newcommand{\CSLLeftMargin}[1]{\parbox[t]{\csllabelwidth}{\strut#1\strut}}
\newcommand{\CSLRightInline}[1]{\parbox[t]{\linewidth - \csllabelwidth}{\strut#1\strut}}
\newcommand{\CSLIndent}[1]{\hspace{\cslhangindent}#1}
\setlength{\emergencystretch}{3em} % prevent overfull lines
\providecommand{\tightlist}{%
  \setlength{\itemsep}{0pt}\setlength{\parskip}{0pt}}
\usepackage{bookmark}
\IfFileExists{xurl.sty}{\usepackage{xurl}}{} % add URL line breaks if available
\urlstyle{same}
\hypersetup{
  pdftitle={Scikit-longitudinals TpT splitter: growing decision trees along the time dimension of longitudinal data},
  pdfauthor={Mathias Valla; Simon Provost},
  hidelinks,
  pdfcreator={LaTeX via pandoc}}

\title{Scikit-longitudinals TpT splitter: growing decision trees along
the time dimension of longitudinal data}
\author{Mathias Valla \and Simon Provost}
\date{2026-04-22}

\begin{document}
\maketitle

\section{Introduction}\label{introduction}

\texttt{scikit-longitudinal} extends the \texttt{scikit-learn} ecosystem
with estimators and workflows tailored to repeated-measures data
(Provost and Freitas 2025). This article focuses on the Time-penalized
Trees (TpT) splitter, a decision-tree split rule designed for
longitudinal covariates (Valla 2024). Rather than reproducing the full
methodological development, we document the current TpT API in the
package and provide two compact, fully reproducible real-data examples.
Beyond the original algorithmic presentation, a recent theoretical
analysis shows that TpT is a consistent estimator along a non-empty set
of observation times under explicit structural assumptions (Valla and
Milhaud 2026).

Tree-based methods for longitudinal or clustered data already span
several families. Model-based recursive partitioning and generalized
linear model trees target subgroup structure through
parameter-instability ideas (Zeileis, Hothorn, and Hornik 2008; Seibold,
Hothorn, and Zeileis 2019). RE-EM trees and mixed-effects random forests
incorporate random effects to account for within-subject or
within-cluster dependence (Sela and Simonoff 2012; Hajjem, Bellavance,
and Larocque 2014). Survival-tree extensions such as LTRC trees can also
accommodate time-varying covariates in censored settings (Fu and
Simonoff 2017). TpT occupies a narrower but complementary niche: it
keeps a CART-style impurity objective and modifies the split-selection
criterion through an explicit time penalty, making the fitted tree
directly readable in terms of both predictor and wave (Valla 2024).

The repository is structured as required by Computo (Computo Team 2020):
the paper is written in Quarto, the environment is frozen in
\texttt{environment.yml}, and the data preparation and analysis code are
kept alongside the manuscript.

\section{\texorpdfstring{The TpT splitter in
\texttt{scikit-longitudinal}}{The TpT splitter in scikit-longitudinal}}\label{the-tpt-splitter-in-scikit-longitudinal}

Recursive partitioning grows a tree by repeatedly splitting a parent
node into child nodes that are more homogeneous with respect to the
target variable. In other words, each split is chosen to reduce
within-node heterogeneity, or equivalently to maximize the loss of
impurity from the parent node to its children (Breiman et al. 1984).

For a classification node \(g\), let \(p_k(g)\) denote the proportion of
class \(k \in \{1, \ldots, K\}\) inside the node. An impurity measure
should be nonnegative, equal to \(0\) when the node is pure, maximal at
the uniform class distribution, and symmetric in the class labels. Two
standard choices are the entropy

\[
I_{\text{Ent}}(g) = -\sum_{k=1}^{K} p_k(g)\log p_k(g)
\]

and the Gini impurity

\[
I_{\text{Gini}}(g) = \sum_{k=1}^{K} p_k(g)\left(1 - p_k(g)\right)
= 1 - \sum_{k=1}^{K} p_k(g)^2.
\]

For regression, CART typically uses the node mean squared error

\[
I_{\text{MSE}}(g) = \frac{1}{N(g)}
\sum_{i: x^{(i)} \in g} \left(y^{(i)} - \bar y_g\right)^2,
\]

where \(\bar y_g\) is the empirical mean response in node \(g\). If a
parent node \(g_p\) is split into left and right child nodes \(g_l\) and
\(g_r\), the usual CART goodness-of-split is the impurity decrease

\[
\Delta I(g_p; g_l, g_r) =
I(g_p) -
\left[
\frac{N(g_l)}{N(g_p)} I(g_l) +
\frac{N(g_r)}{N(g_p)} I(g_r)
\right],
\]

with \(I(g)\) standing for one of the classification impurities above,
or for the MSE in regression.

At a tree node associated with parent time index \(t_p\), TpT modifies
this criterion and scores a candidate split at time \(t_c\) by
multiplying the usual impurity improvement \(\Delta I\) by an
exponential time penalty:

\[
G_{\gamma} = \Delta I \exp \left[-\gamma (t_c - t_p)\right].
\]

The hyperparameter \(\gamma \ge 0\) controls how strongly later
measurements are penalized. When \(\gamma = 0\), TpT reduces to ordinary
impurity-based split selection. As \(\gamma\) increases, the tree
becomes more reluctant to use later waves unless they bring a clearly
stronger impurity reduction.

From a user perspective, the main entry points are
\texttt{TpTDecisionTreeClassifier} and
\texttt{TpTDecisionTreeRegressor}. Both can consume LONG-format data
directly by specifying an identifier column, an observation-time column
and a subject-specific duration column.

The same splitter can also be reached from an explicit wide matrix
together with \texttt{features\_group} metadata. We use both entry
points below: the \texttt{pbcseq} example illustrates direct long-format
fitting, while the \texttt{Produc} example makes the long-to-wide
conversion explicit.

For this article, we select \(\gamma\) on a small candidate grid by
minimizing the fitted tree's weighted total leaf impurity on the
training split. This is a pragmatic paper-specific heuristic whose only
goal is to keep the examples deterministic and easy to rerun; it should
not yet be interpreted as a recommended general tuning rule for applied
work.

\subsection{A compact theoretical
guarantee}\label{a-compact-theoretical-guarantee}

TpT is not only a practical heuristic. In Valla and Milhaud (2026), the
analysis focuses on record-high times, that is, feasible times whose
penalized population envelope dominates all later competitors. The main
result requires separated penalized-envelope gaps (Assumption A1),
finite-sample attainability of the corresponding envelope peaks
(Assumption A2), the TpT-\(\alpha\) branch-growing rule, and the
regularity and noise assumptions needed to invoke the Klusowski-Tian
oracle inequality for depth-\(K\) CART estimators (Klusowski and Tian
2024).

\textbf{Theorem 1 Valla and Milhaud (2026).} Assume (A1)--(A2), the
TpT-\(\alpha\) growing rule, and the regularity and noise conditions
required to apply the Klusowski-Tian oracle inequality for depth-\(K\)
CART estimators. For every fixed record-high time \(t_k\) such that the
approximation bias vanishes as \(K(t_k) \to \infty\), the estimator
fitted on the realized \(t_k\)-subtree satisfies

\[
\mathbb{E}\!\left[\left\lVert \hat f_{t_k, K(t_k)} - f_{t_k}^{\star} \right\rVert^2_{L^2(P_{X(t_k)})}\right]
\longrightarrow 0
\qquad \text{as } n \to \infty.
\]

This is intentionally narrower than a blanket consistency claim over all
time indices. The result identifies a non-empty sequence of times at
which the penalized TpT rule recovers a consistent subtree estimator. We
include this result only to summarize the current theoretical status of
TpT; the paper itself remains focused on software usage. For an
implementation-focused article, the practical takeaway is that the TpT
framework underlying the splitter now has a precise asymptotic analysis,
even though the present package implementation is discussed here from an
operational rather than theorem-verification standpoint.

\subsection{Current scope}\label{current-scope}

The present implementation is suitable for compact reproducible examples
and software-oriented demonstrations. Missing-value routing inside the
Cython TpT splitter is still treated as future work, and monotonic
constraints are explicitly disabled for TpT at the wrapper level.

\section{Reproducible setup}\label{reproducible-setup}

The repository contains:

\begin{itemize}
\tightlist
\item
  \texttt{scripts/prepare\_example\_data.py} to download and preprocess
  the open datasets;
\item
  \texttt{analysis/tpt\_examples.py} with helper functions used by the
  paper;
\item
  \texttt{environment.yml} with a frozen Python environment.
\end{itemize}

The quick-start installation shown below is the user-facing entry point.
Exact reproduction of this paper instead relies on
\texttt{environment.yml}, which fixes Python 3.11, the released
\texttt{scikit-lexicographical-trees==0.0.5}, and the
\texttt{scikit-longitudinal} commit used for the TpT wrapper/API state
documented here. The helper script \texttt{setup-render-ci.sh} ensures
that Quarto renders against the extended tree backend required by the
TpT examples.

\section{Install, quick start, and main
parameters}\label{install-quick-start-and-main-parameters}

The official \texttt{scikit-longitudinal} getting-started guide is
intentionally lightweight, and the same philosophy fits this article
well. For a standard Python environment, the minimal installation is:

\begin{Shaded}
\begin{Highlighting}[]
\ExtensionTok{pip}\NormalTok{ install scikit{-}longitudinal}
\end{Highlighting}
\end{Shaded}

For LONG-format data, the shortest end-to-end TpT workflow is to pass a
subject identifier, an observation-time column, a subject-level duration
column, and the names of the longitudinal predictors:

The small set of parameters that matters most in practice is summarized
in (\textbf{tbl-tpt-main-parameters?}).

\phantomsection\label{tbl-tpt-main-parameters}

When the data are already stored in wave-separated columns, or when TpT
is used inside a broader \texttt{scikit-longitudinal} workflow, the
package documentation recommends the same building blocks as the rest of
the library: a \texttt{LongitudinalPipeline}, a time-aware preprocessing
step such as \texttt{MerWavTimePlus}, and a final estimator. A minimal
pipeline skeleton is:

This article keeps the examples at the estimator level, but the API is
compatible with these larger package workflows. In particular, the two
examples below are meant to bridge the package documentation and
practical usage: \texttt{pbcseq} uses the long-format wrapper directly,
whereas \texttt{Produc} shows the explicit conversion and the low-level
\texttt{features\_group} interface.

\section{Data sets}\label{data-sets}

The two examples use small public repeated-measures datasets with
complementary temporal structures, summarized in
(\textbf{tbl-dataset-summary?}). We cite \texttt{scikit-longitudinal}
through its JOSS software paper (Provost and Freitas 2025). The
\texttt{pbcseq} data contain follow-up laboratory measurements for the
randomized Mayo Clinic primary biliary cirrhosis trial, with visits
planned at 6 months, 1 year, and yearly thereafter (Therneau, Terry
2025). The \texttt{Produc} data form a balanced annual panel of 48 US
states from 1970 to 1986 (Munnell, Alicia H. 2025). Together they
provide two deliberately different views of TpT: a subject-level
clinical classification problem with irregular follow-up, and a compact
multi-wave regression panel where the effect of \(\gamma\) on selected
split times is easy to inspect.

\phantomsection\label{tbl-dataset-summary}

\section{\texorpdfstring{Example 1: Subject-level classification on
\texttt{pbcseq}}{Example 1: Subject-level classification on pbcseq}}\label{example-1-subject-level-classification-on-pbcseq}

Our classification example uses \texttt{pbcseq}, the sequential primary
biliary cirrhosis follow-up data distributed with the \texttt{survival}
package (Therneau, Terry 2025). We keep four complete longitudinal
biomarkers (\texttt{bili}, \texttt{albumin}, \texttt{protime},
\texttt{ast}), use \texttt{day} as the observation-time variable,
\texttt{futime} as the subject-specific horizon, and define a binary
target equal to one when the subject experiences either transplantation
or death by the end of follow-up. In the processed version used here,
this yields 312 subjects, 1,945 recorded visits, and a median of 5
visits per subject ((\textbf{tbl-dataset-summary?})).

The example is run on a subject-level train/test split, with the
long-format API of \texttt{TpTDecisionTreeClassifier}. The minimal
fit/predict workflow is:

\phantomsection\label{tbl-pbcseq-metrics}

We also report the small hyperparameter set used for this preliminary
example.

\phantomsection\label{tbl-pbcseq-hyperparameters}

The selected gamma comes from the training-grid search shown in
(\textbf{tbl-pbcseq-gamma-search?}) and visualized in
(\textbf{fig-pbcseq-gamma-profile?}).

\phantomsection\label{tbl-pbcseq-gamma-search}

\begin{figure}
\centering
\pandocbounded{\includegraphics[keepaspectratio]{paper_files/figure-pdf/fig-pbcseq-gamma-profile-output-1.pdf}}
\caption{Lower-envelope training tree-total-impurity profile for the
pbcseq classification example, extracted from the full 0.001 gamma grid
by keeping the local minima of the oscillating
path.}\label{fig-pbcseq-gamma-profile}
\end{figure}

The resulting subject-level metrics are shown in
(\textbf{tbl-pbcseq-metrics?}). A small preview of the held-out
predictions is given in (\textbf{tbl-pbcseq-predictions?}). With this
finer gamma grid, the same depth-3, \texttt{min\_samples\_leaf=10}
configuration selects \texttt{gamma\ =\ 0.184} and still uses later
history: the learned tree splits at waves 2, 3, 4, and 7 rather than
collapsing to baseline-only rules. This is the behavior we want to
expose in the article: the time penalty regularizes the use of history
without erasing it when later measurements still provide a clearly
better impurity reduction.

\phantomsection\label{tbl-pbcseq-predictions}

The fitted tree is displayed below with the dedicated
\texttt{plot\_tree} exporter, using the wide feature names built
internally by the TpT preprocessing step.

\clearpage

\begin{figure}
\centering
\pandocbounded{\includegraphics[keepaspectratio]{paper_files/figure-pdf/fig-pbcseq-tree-output-1.pdf}}
\caption{Selected TpT classification tree on pbcseq, displayed with the
dedicated plot\_tree exporter.}\label{fig-pbcseq-tree}
\end{figure}

This is precisely the kind of output we want in this article: the fitted
tree shows explicitly which wave-specific biomarkers were selected by
TpT, while the calling interface remains close to ordinary
\texttt{scikit-learn} usage. The node labels make the classifier
interpretable in both dimensions at once: which biomarker matters, and
at which follow-up wave it becomes decisive for the split. In this run,
the selected nonzero gamma still preserves several later visits:
bilirubin enters at waves 2 and 4, albumin at waves 3 and 4, prothrombin
time at wave 7, and AST at wave 3.

\section{\texorpdfstring{Example 2: State-level regression on
\texttt{Produc}}{Example 2: State-level regression on Produc}}\label{example-2-state-level-regression-on-produc}

Our regression example uses \texttt{Produc}, a public US state
production panel available through Rdatasets (Munnell, Alicia H. 2025).
The source data contain 48 states followed yearly from 1970 to 1986. We
keep three longitudinal predictors available at each wave: private
capital stock (\texttt{pc}), employment (\texttt{emp}), and unemployment
rate (\texttt{unemp}). For each state, we hold out the final recorded
gross state product and predict that \texttt{final\_gsp} from the
preceding 16 yearly observations. We keep \texttt{Produc} deliberately:
although it is an economic panel rather than an individual-level
biomedical cohort, its 17 annual waves and three interpretable
predictors make the effect of the TpT time penalty easy to see in a
small real dataset. To keep the example readable while still producing a
nonzero selected gamma and non-baseline splits on the training impurity
criterion, we use the predictor subset \texttt{pc}, \texttt{emp}, and
\texttt{unemp} together with a depth-4 tree,
\texttt{min\_samples\_leaf=1}, and \texttt{min\_samples\_split=2}.

In this example we explicitly use the TpT conversion helper
\texttt{long\_to\_wide(...)}, then fit the regressor on the resulting
wide matrix and its \texttt{feature\_groups} metadata. The minimal
workflow is:

The next two tables show the same training subjects before and after the
explicit long-to-wide conversion.

\phantomsection\label{tbl-produc-long-preview}

\phantomsection\label{tbl-produc-wide-preview}

\phantomsection\label{tbl-produc-metrics}

\phantomsection\label{tbl-produc-hyperparameters}

\phantomsection\label{tbl-produc-gamma-search}

\begin{figure}
\centering
\pandocbounded{\includegraphics[keepaspectratio]{paper_files/figure-pdf/fig-produc-gamma-profile-output-1.pdf}}
\caption{Training tree-total-impurity profile (top) and selected split
times (bottom) for the Produc regression example on the article grid
from 0.000 to 0.100 by 0.001.}\label{fig-produc-gamma-profile}
\end{figure}

\phantomsection\label{tbl-produc-predictions}

\begin{figure}
\centering
\pandocbounded{\includegraphics[keepaspectratio]{paper_files/figure-pdf/fig-produc-predictions-output-1.pdf}}
\caption{Observed versus predicted final gross state product in the
Produc regression example, using the full-history
inputs.}\label{fig-produc-predictions}
\end{figure}

The selected regression tree can also be shown directly with the
dedicated \texttt{plot\_tree} exporter.

On the article grid from \texttt{0.000} to \texttt{0.100} by
\texttt{0.001}, the selected regression gamma is \texttt{0.011}. The
corresponding tree still uses later history, with selected split times
between waves 1 and 5. As gamma increases further on that grid, the tree
is pulled back toward baseline: the root split moves to wave 0 by about
\texttt{gamma\ =\ 0.031}, and from about \texttt{gamma\ =\ 0.055} onward
all selected split times are at baseline. The gamma path therefore shows
the intended TpT tradeoff quite clearly: modest penalization still
allows mid-history splits when they materially reduce impurity, whereas
stronger penalization progressively collapses the tree toward earlier
waves.

\begin{figure}
\centering
\pandocbounded{\includegraphics[keepaspectratio]{paper_files/figure-pdf/fig-produc-tree-output-1.pdf}}
\caption{Selected TpT regression tree on Produc, displayed with the
dedicated plot\_tree exporter.}\label{fig-produc-tree}
\end{figure}

This makes the regression example a useful complement to
\texttt{pbcseq}: the exported tree tells us not only which economic
predictor matters, but also whether it matters near the beginning of the
panel or later in the observed trajectory. In the selected tree, the
root decision already depends on employment at wave 1, while deeper
refinements still use later information through private capital at wave
2 and employment again at wave 5.

\section{Discussion}\label{discussion}

These two examples are intentionally small. Their role is not to provide
a full benchmarking study, but to document a reproducible and
inspectable usage path for TpT within \texttt{scikit-longitudinal}.
Together they show that TpT can be fit either directly from LONG-format
data or from an explicit wide matrix, tuned with a simple deterministic
gamma heuristic for demonstration purposes, and exported as readable
wave-indexed trees.

Most importantly for this article, the exported trees show both what
matters and when it matters. In \texttt{pbcseq}, the classifier can
isolate clinically meaningful biomarker thresholds at specific follow-up
waves rather than collapsing the history into a single summary. In
\texttt{Produc}, the regression tree reveals whether private capital,
employment, or unemployment becomes informative early or later in the
state-level trajectory. This predictor-plus-wave readability is the main
practical property we wanted to make visible.

The paper also has clear limits. We deliberately keep the methodological
discussion short, we do not attempt a large comparative benchmark, and
the gamma-selection rule used here is only a reproducible placeholder
rather than a definitive model-selection recommendation. On the software
side, missing-value routing inside the Cython TpT splitter is still
future work.

Natural next steps are equally clear: implement survival-specific
evaluation metrics, add richer visualization utilities such as
subject-level trajectory displays inside the fitted tree, build ensemble
methods on top of the current splitter, and generalize the time penalty
so that prior temporal structure such as lagged predictor effects can be
encoded directly. Because TpT is now integrated into the tree-building
backend, these extensions are largely a question of design and
validation rather than reimplementing the core split rule from scratch.

Within that scope, TpT already has the ingredients of a useful
\texttt{scikit-longitudinal} contribution: a lightweight API,
interpretable exported trees, and a growing theoretical foundation.

\section{Acknowledgements}\label{acknowledgements}

Work(s) conducted within the Research Chair DIALog under the aegis of
the Risk Foundation, a joint initiative by Université Paris-Dauphine PSL
and CNP Assurances.

Additional acknowledgements for Simon Provost can be added in the
submission-ready version.

\section*{References}\label{references}
\addcontentsline{toc}{section}{References}

\phantomsection\label{refs}
\begin{CSLReferences}{1}{0}
\bibitem[\citeproctext]{ref-breiman1984cart}
Breiman, Leo, Jerome H. Friedman, Richard A. Olshen, and Charles J.
Stone. 1984. \emph{Classification and Regression Trees}. Wadsworth
International Group.

\bibitem[\citeproctext]{ref-computo}
Computo Team. 2020. {``Computo: Reproducible Computational/Algorithmic
Contributions in Statistics and Machine Learning.''}

\bibitem[\citeproctext]{ref-fu2017ltrc}
Fu, Wei, and Jeffrey S. Simonoff. 2017. {``Survival Trees for
Left-Truncated and Right-Censored Data, with Application to Time-Varying
Covariate Data.''} \emph{Biostatistics} 18 (2): 352--69.
\url{https://doi.org/10.1093/biostatistics/kxw047}.

\bibitem[\citeproctext]{ref-hajjem2014merf}
Hajjem, Ahlem, François Bellavance, and Denis Larocque. 2014.
{``Mixed-Effects Random Forest for Clustered Data.''} \emph{Journal of
Statistical Computation and Simulation} 84 (6): 1313--28.
\url{https://doi.org/10.1080/00949655.2012.741599}.

\bibitem[\citeproctext]{ref-klusowski2024cart}
Klusowski, Jason M., and Peter M. Tian. 2024. {``Large Scale Prediction
with Decision Trees.''} \emph{Journal of the American Statistical
Association} 119 (545): 525--37.
\url{https://doi.org/10.1080/01621459.2022.2126782}.

\bibitem[\citeproctext]{ref-dataset-produc}
Munnell, Alicia H. 2025. {``Produc: US State Production Panel.''}
\url{https://vincentarelbundock.github.io/Rdatasets/doc/Ecdat/Produc.html}.

\bibitem[\citeproctext]{ref-provost2025sklong}
Provost, Simon, and Alex A. Freitas. 2025. {``Scikit-Longitudinal: A
Machine Learning Library for Longitudinal Classification in Python.''}
\emph{Journal of Open Source Software} 10 (112): 8481.
\url{https://doi.org/10.21105/joss.08481}.

\bibitem[\citeproctext]{ref-seibold2019glmtrees}
Seibold, Heidi, Torsten Hothorn, and Achim Zeileis. 2019. {``Generalised
Linear Model Trees with Global Additive Effects.''} \emph{Advances in
Data Analysis and Classification} 13 (3): 703--25.
\url{https://doi.org/10.1007/s11634-018-0342-1}.

\bibitem[\citeproctext]{ref-sela2012reem}
Sela, Rebecca J., and Jeffrey S. Simonoff. 2012. {``{RE-EM} Trees: A
Data Mining Approach for Longitudinal and Clustered Data.''}
\emph{Machine Learning} 86 (2): 169--207.
\url{https://doi.org/10.1007/s10994-011-5258-3}.

\bibitem[\citeproctext]{ref-dataset-pbcseq}
Therneau, Terry. 2025. {``Pbcseq: Mayo Clinic Primary Biliary Cirrhosis
Sequential Data.''} \url{https://rdrr.io/cran/survival/man/pbcseq.html}.

\bibitem[\citeproctext]{ref-valla2024tpt}
Valla, Mathias. 2024. {``Time-Penalized Trees ({TpT}): Introducing a New
Tree-Based Data Mining Algorithm for Time-Varying Covariates.''}
\emph{Annals of Mathematics and Artificial Intelligence} 92: 1609--61.
\url{https://doi.org/10.1007/s10472-024-09950-w}.

\bibitem[\citeproctext]{ref-valla2026consistency}
Valla, Mathias, and Xavier Milhaud. 2026. {``Consistent Time-Aware Trees
for Longitudinal Data: The Time-Penalized Tree.''}
\url{https://cnrs.hal.science/hal-05022929}.

\bibitem[\citeproctext]{ref-zeileis2008mob}
Zeileis, Achim, Torsten Hothorn, and Kurt Hornik. 2008. {``Model-Based
Recursive Partitioning.''} \emph{Journal of Computational and Graphical
Statistics} 17 (2): 492--514.
\url{https://doi.org/10.1198/106186008X319331}.

\end{CSLReferences}

\end{document}
